\documentclass[10pt]{article}
\usepackage[margin=0.79in]{geometry}
\usepackage[T1]{fontenc}
\usepackage{amsmath,amssymb,booktabs,longtable,hyperref,tikz}
\usetikzlibrary{arrows.meta,positioning}
\hypersetup{colorlinks=true,urlcolor=blue}
\title{\textbf{Fractal Holonet: Five Engineering Frontiers}\\
Four-MZI PIC handoff, coherent interferometry, $ \mathbb{Z}_4 $ FPGA control, loss-aware routing, and qutrit separation}
\author{W33-Theory engineering programme}
\date{23 October 2026 --- executable simulation / hardware-source edition}
\begin{document}\maketitle
\begin{abstract}
This edition extends the 40-line-bus W33 photonic construction by implementing
five connected technical investigations: a technology-neutral foundry I/O
contract, a four-phase coherent route interferometer protocol, a
synthesizable SystemVerilog modulo-four phase-frame controller, a
loss/phase-aware Dijkstra router with line-bus failure injection,
and a loss/leakage-aware qutrit Fourier-gate bridge. It also
derives an alternative fixed Walsh--Hadamard optical architecture
that simultaneously diagonalizes the four XOR bus permutations.
Mathematical and software claims are distinguished from missing
fabrication, simulator, and experimental evidence.
\end{abstract}
\tableofcontents
\section{Starting point and source boundaries}
The preceding files
\texttt{papers/fractal\_holonet\_40bus\_optical\_addendum\_20261023.tex}
and \texttt{papers/fractal\_holonet\_xor2\_photonic\_phase\_frame\_20261023.tex}
establish forty four-port line buses, three switch states and an ideal
four-MZI-per-bus compilation. The original
\texttt{holonet\_machine\_blueprint.pdf},
\texttt{photonic\_holonet.pdf}, and
\texttt{papers/forty\_points/main.pdf} motivate the proposed machine.
Earlier project work in \texttt{analysis/w33\_scheduler\_os.py}
already discovered the 40 buses; that combinatorics is not new here.
The source checkout is shared with unrelated parallel research. The
new outputs are isolated to newly named files.

\section{Track 1: fail-closed PIC / PDK engineering handoff}
The 40-bus logical parent has 160 abstract MZI instances,
80 paired stage-control words and 160 provisional individual
electrode device identifiers. Each instance now has a named
optical input/output pair, required mode pairing, command word
and unresolved process technology fields in
\texttt{data/w33\_fractal\_all5\_engineering\_20261023.json}.
\emph{This is an implementable schematic contract, not a finished GDS.}
The connected workstation did not have gdsfactory, KLayout,
an authorized foundry PDK or a DRC/LVS deck. The handoff
therefore keeps \texttt{foundry\_pdk\_selected},
\texttt{DRC\_passed}, \texttt{optical\_LVS\_passed}, and
\texttt{foundry\_tapeout\_ready} explicitly false. In
particular, second-stage mode pairs $(0,2),(1,3)$
require real topology-specific crossing/routing solutions.

\section{Track 2: falsifiable coherent interferometry}
For ideal switch masks $m$, the bus transfer is
\[
U_m=i^{w(m)}P_m,\qquad m\in\{0,1,2,3\}.
\]
The two W33 length-two routes
$0\rightarrow4\rightarrow7$ and $0\rightarrow14\rightarrow7$
differ by an ideal factor of $i$ in the frozen mode-port
convention. A balanced recombiner with reference phase
$\theta$ and interference visibility $V$ would ideally
produce a normalized single-port detection probability
\[
p_+(\theta)=\frac{1+V\cos(\Delta+\theta)}2,
\qquad \Delta=\frac{\pi}{2}.
\]
The software prepares four phase settings
$\theta\in\{0,\pi/2,\pi,3\pi/2\}$ and samples independent
binomial counts. Quadrature differences estimate
$V\cos\Delta$ and $V\sin\Delta$; using
\[
x=\hat p_+(0)-\hat p_+(\pi),\quad
y=\hat p_+(3\pi/2)-\hat p_+(\pi/2),
\quad \hat\Delta=\operatorname{atan2}(y,x)
\]
correctly reconstructs the sign convention. Its sample data
are simulated and seeded, \emph{not laboratory photon counts}.
The protocol requires stable source/detector calibration,
path equalization, a phase reference, balanced recombination
and a physically controllable delay or phase shifter.
Observed phase alone does not prove exotic temporal
self-entanglement or a new topological holonomy.

\section{Track 3: coherent $\mathbb{Z}_4$ phase-frame controller}
The new controller
\texttt{hardware/w33\_linebus\_phase\_frame.sv}
exposes a valid/ready registered command interface with
\texttt{cmd\_mask[1:0]}, \texttt{cmd\_clear},
\texttt{result\_frame[1:0]}, and
\texttt{result\_correction[1:0]}. If a command is
accepted, the state transition is
\[
e_{\mathrm{next}}
=(e_{\mathrm{old}}+w(m))\bmod4,
\qquad c_{\mathrm{next}}=-e_{\mathrm{next}}\bmod4.
\]
A clear command resets the accumulation before adding the
current switch's phase. Backpressure holds all registered
outputs stable until consumed. The ready/valid semantics are
also implemented in a Python reference with deterministic and
randomized tests. A self-checking SystemVerilog testbench is
provided at \texttt{hardware/tb\_w33\_linebus\_phase\_frame.sv}.
\textbf{No Verilator, Icarus or Yosys tool was installed}; therefore
this release does not claim compiled HDL, post-synthesis timing
or functional HDL simulation. A phase-frame word does not
automatically correct a physical optical path: coherent
relative phases must be referenced or compensated where relevant.

\section{Track 4: loss-, phase- and failure-aware routing}
The explicit line-bus graph has 1600 leaves at recursion
depth two and exact diameter six. The new optimizer
treats each search state as
\[
(\text{leaf address}, e),\qquad e\in\mathbb{Z}_4
\]
and uses Dijkstra's algorithm on a synthetic
nonnegative link-loss cost. An optional endpoint
condition $e=0$ forces routes with the desired
modulo-four coherent phase. A failed line-bus ID
removes all logical links through that line
and exercises alternate-route selection.
Costs include declared per-MZI attenuation,
illustrative line-/mask-dependent penalties,
and an explicitly optional time-slot weight.
The computed attenuation is neither a measured
physical path loss nor a quantum channel fidelity.
Because an exact phase-frame constraint can force
a longer path, route-count comparisons are reported
separately from modeled power transmission.
No guarantee is made for congestion, multiuser
scheduling, coherent multi-path splitting, or
fault-tolerant live quantum computation.

\section{Track 5: typed qutrit Fourier-gate bridge}
On three computational optical modes, set
\[
F_3=\frac1{\sqrt3}
\begin{pmatrix}
1&1&1\\1&\omega&\omega^2\\
1&\omega^2&\omega
\end{pmatrix},\qquad
\omega=e^{2\pi i/3}.
\]
The separate four-mode qutrit-gate target is
\[
U_{F_3}=\operatorname{diag}(F_3,1).
\]
An independent complex Givens elimination
constructs this target with exactly three
two-mode SU(2) blocks acting \emph{only} on
modes 0, 1 and 2, plus a phase diagonal.
Numerical reconstruction is checked.
This is an abstract optical interferometer
program, not the already verified four-state
XOR bus switch.

A strong negative result prevents a tempting
architectural error. For any nonidentity XOR
permutation $P_m$ on four computational basis
modes, its $3\times3$ upper-left block has
singular values $(1,1,0)$. Thus among input
basis states $|0\rangle,|1\rangle,|2\rangle$,
one is routed into the fourth auxiliary port.
The success probability for that basis input
is zero, and the uniform-basis mean is $2/3$.
Consequently \emph{the binary four-port router
cannot act as a deterministic qutrit gate
on the naive three-spatial-mode encoding}.
The correct solution requires a distinct
three-mode programmable gate or an explicit
different encoding and physical interface.

\section{Outside-box result: Walsh optics instead of active MZIs}
Let $H_2=2^{-1/2}
\left(\begin{smallmatrix}1&1\\1&-1\end{smallmatrix}\right)$
and $W=H_2\otimes H_2$. Since the four XOR
permutations commute, all are simultaneously
diagonalized:
\[
P_m=W\operatorname{diag}
\left((-1)^{s\cdot m}\right)_{s\in\mathbb{F}_2^2}W.
\]
The exact switch family can therefore be expressed
with \emph{fixed} input/output four-port Walsh
mixers and a diagonal set of two independent
binary phase words. A conventional
two-stage implementation of each Walsh mixer
uses four passive 50:50 couplers, giving eight
fixed couplers over the pair plus diagonal
phase elements. This trades four \emph{active}
MZIs for more \emph{passive} couplers and two
binary phase settings, without establishing
a better real device. Which choice minimizes
energy, footprint, loss, and complexity is a
process-specific experimental question.

\section{Next physical and proof obligations}
\begin{longtable}{p{0.29\linewidth}p{0.54\linewidth}}
\toprule
Area & Evidence required\\\midrule
Foundry & Authorized PDK, actual optics S-parameters,
real GDS/OASIS, optical LVS/DRC and electrical parasitics\\
Quantum test & Interferometric reference, source/measurement
SPAM model, visibility and delay scan, real detector counts\\
HDL & Independent Icarus/Verilator simulation,
synthesis, clock constraints, resource and timing closure\\
Routing & Traffic, optical coherence, congestion,
physical loss, realistic switch time and fault coverage\\
Qutrit & Optical F3 cells, finite loss/leakage,
auxiliary-mode port isolation and single-photon test\\
\bottomrule
\end{longtable}
\section{Files and reproduction}
\texttt{analysis/w33\_fractal\_all5\_engineering\_20261023.py}
recomputes the mathematical/synthetic tracks and exports
\texttt{data/w33\_fractal\_all5\_engineering\_20261023.json}.
The regression file is
\texttt{tests/test\_w33\_fractal\_all5\_engineering\_20261023.py}.
The separate reference model for hardware timing
is \texttt{tests/test\_w33\_linebus\_phase\_frame\_protocol\_20261023.py}.
The lab, foundry and physical-hardware gates remain open.
\end{document}